
\section{常值映射模空间的定向}
\label{sec:8-6}

下面，设 $P_1$ 和 $P_2$ 是定向子流形。考虑常值映射模空间%
\liftIndex{I303}
$\mathcal M_3^{\mathrm{main}}(L,\beta_0;P_1,P_2)$，作为集合有
\[
 \mathcal M_3^{\mathrm{main}}(L,\beta_0;P_1,P_2)=P_1\cap P_2.
\]
我们需要研究两边定向之间的差异。为此，先确认 $P_1\cap P_2$
上的定向约定。

\begin{statement}{约定}{8.6.1}\label{conv:8-6-1}
设 $P_1$ 和 $P_2$ 是 $L$ 中横截相交的子流形。分别用
$n,p_1,p_2$ 表示 $L,P_1,P_2$ 的维数，并假定它们均已定向。
取 $x\in P_1\cap P_2$。$P_i$ 在 $L$ 中的法丛 $N_{P_i}L$
（$i=1,2$）之定向由
\[
 T_xL=N_{x,P_i}L\times T_xP_i
\]
定义。由于 $P_1$ 和 $P_2$ 横截相交，可假定
$N_{x,P_1}L\subset T_xP_2$ 且 $N_{x,P_2}L\subset T_xP_1$。
然后由
\[
 T_xL=N_{x,P_1}L\times N_{x,P_2}L\times T_x(P_1\cap P_2)
\]
定义 $P_1\cap P_2$ 的定向。
\end{statement}

\begin{remark}{8.6.2}\label{rem:8-6-2}
上述约定在 $P_1\cap P_2$ 上定义了“上同调定向”。也就是说，
\[
 P_1\cap P_2=(-1)^{(n-p_1)(n-p_2)}P_2\cap P_1.
\]
\end{remark}

现在可以证明下列结论。

\begin{proposition}{8.6.3}\label{prop:8-6-3}
可以在模空间 $\mathcal M_3^{\mathrm{main}}(L;\beta_0)$ 上给出一个
定向仓西结构，使下列同构保定向：
\[
 \operatorname{ev}_0\colon
 \mathcal M_3^{\mathrm{main}}(L;\beta_0)_{
  (\operatorname{ev}_1,\operatorname{ev}_2)}
 \mathop{\times}_{f_1\times f_2}(P_1\times P_2)
 \longrightarrow
 (-1)^{n(p_1+p_2)+p_1(n-p_2)}P_1\cap P_2.
\]
\end{proposition}

注意，由定义~\ref{def:8-4-1}，
\[
\begin{aligned}
 \mathcal M_3^{\mathrm{main}}(L;\beta_0;P_1,P_2)
 ={}&(-1)^{(n+1)\deg P_1}
 \mathcal M_3^{\mathrm{main}}(L;\beta_0)_{
  (\operatorname{ev}_1,\operatorname{ev}_2)}\\[-2pt]
 &{}\mathop{\times}_{f_1\times f_2}(P_1\times P_2),
\end{aligned}
\]
因此立即得到下列结果。

\begin{corollary}{8.6.4}\label{cor:8-6-4}
在 $\mathcal M_3^{\mathrm{main}}(L;\beta_0)$ 上使用命题
\ref{prop:8-6-3} 的定向仓西结构，有
\[
 \bigl(\mathcal M_3^{\mathrm{main}}(L;\beta_0;P_1,P_2),
 \operatorname{ev}_0\bigr)
 =(-1)^{\deg P_1(\deg P_2+1)}P_1\cap P_2.
\]
\end{corollary}

\begin{proof}[命题~\ref{prop:8-6-3} 的证明]
为了证明命题~\ref{prop:8-6-3}，我们将显式描述
$\mathcal M_3^{\mathrm{main}}(L;\beta_0)$、$P_1\times P_2$ 上（弱浸没的）
仓西结构及它们的纤维积。

首先，显式描述 $\mathcal M_3^{\mathrm{main}}(L;\beta_0)$ 上的仓西结构，
并研究其定向。由于
$\mathcal M_3^{\mathrm{main}}(L;\beta_0)$ 由常值映射组成，这是一个局部问题。
因而只须考虑 $L=\mathbb R^n$、$M=\mathbb C^n$ 且两者均带标准定向的
情形。在单位圆盘 $D^2$ 的边界上取三个标记点
\[
 z_0=1,\qquad z_1=\sqrt{-1},\qquad z_2=-1.
\]
$\mathcal M_3^{\mathrm{main}}(L;\beta_0)$ 上的
$\operatorname{ev}_1$ 和 $\operatorname{ev}_2$ 分别是在 $z_1$ 和 $z_2$ 处的
评价映射。

置
\[
\begin{aligned}
 U&:=\left\{u_{a,b}\colon D^2\to\mathbb C^n\ \middle|\
 u_{a,b}(z)=\frac12a\sqrt{-1}\,\bar z+b
 -\frac12a\sqrt{-1}\,z,
 \ a,b\in\mathbb R^n\right\},\\
 \mathcal E&:=\left\{(u_{a,b},c)\ \middle|\
 a,b,c\in\mathbb R^n\right\}.
\end{aligned}
\]
并定义截面 $s\colon U\to\mathcal E$ 为
$s(u_{a,b})=(u_{a,b},a)$。%
\translatornote{源书在 $\mathcal E$ 的集合条件中另写
$z\in D^2$，但元组 $(u_{a,b},c)$ 中没有这个自由变量，后文又明定
$\mathcal E$ 是 $U$ 上纤维为 $\mathbb R^n$ 的平凡向量丛；故删去该无关条件。}
于是，以 $L$ 为值域的常值映射模空间就是
\[
 \mathcal M_3^{\mathrm{main}}(L;\beta_0)=s^{-1}(0)
 \cong\mathbb R_b^n.
\]

当然，$U$ 可以识别为 $\mathbb R_a^n\times\mathbb R_b^n$，并可将其
看成 $T_*L\times T_*L$；$\mathcal E$ 是 $U$ 上纤维为
$\mathbb R^n\cong T_*L$ 的平凡向量丛。将第~\ref{subsec:8-1-1}
小节的论证应用于 $\mathcal M_3^{\mathrm{main}}(L;\beta_0)$：藉由
$\operatorname{ev}_0$ 把该模空间与目标 $L$ 识别后，这里的定向就是
$L$ 的定向。我们必须分别给出 $U$ 的定向与 $\mathcal E$ 的逐纤维
定向，使 $\mathcal M_3^{\mathrm{main}}(L;\beta_0)$ 的定向与仓西结构
$(s;\mathcal E\to U)$ 所定义的定向相容。从现在起，通过识别
$U\cong T_*L\times T_*L$ 给 $U$ 定向，通过识别
$\mathbb R^n\cong T_*L$ 给 $\mathcal E$ 的纤维定向。于是
$(s;\mathcal E\to U)$ 给出 $\mathcal M_3^{\mathrm{main}}(L;\beta_0)$ 上的
仓西结构。

注意，评价映射
$(\operatorname{ev}_1,\operatorname{ev}_2)\colon U\to L\times L$ 由
\[
 (\operatorname{ev}_1,\operatorname{ev}_2)(u_{a,b})=(a+b,b)
\]
给出，它保定向。我们有保定向的识别
\[
 \mathcal E_*\times T_*\mathcal M_3^{\mathrm{main}}(L;\beta_0)
 =T_*U=T_*L\times T_*L.
\]

其次，描述 $P_1\times P_2$ 上的仓西结构。设
$U_i\subset L$ 是 $P_i$ 在 $L$ 中的管状邻域（$i=1,2$），设
$N_{P_i}$ 是 $P_i$ 在 $L$ 中的法丛，而
$s_i\colon U_i\to N_{P_i}$ 是典范截面。则
$(s_i;N_{P_i}\to U_i)$ 给出 $P_i$ 上的仓西结构。换言之，
$P_i=s_i^{-1}(0)$，而它的定向由
\[
 N_{P_i}|_{x_i}\times T_{x_i}s_i^{-1}(0)
 =T_{x_i}U_i\cong T_{x_i}L
\]
给出，其中 $x_i\in P_i$。因此，$P_i=s_i^{-1}(0)$ 是定向空间。
仓西结构的乘积
\[
 (s_1\oplus s_2;N_{P_1}\oplus N_{P_2}\longrightarrow U_1\times U_2)
\]
给出 $P_1\times P_2$ 上的仓西结构。依约定~\ref{conv:8-2-1}(4)，
由仓西结构定义的定向与 $P_1\times P_2$ 上乘积定向的关系为
\[
 P_1\times P_2=(-1)^{p_1(n-p_2)}(s_1\oplus s_2)^{-1}(0),
\]
因为
\[
\begin{aligned}
 &N_{P_1}|_{x_1}\times N_{P_2}|_{x_2}
 \times T(s_1\oplus s_2)^{-1}(0)_{(x_1,x_2)}\\
 &\quad=T_{x_1}U_1\times T_{x_2}U_2\\
 &\quad=(N_{P_1}|_{x_1}\times T_{x_1}P_1)
 \times(N_{P_2}|_{x_2}\times T_{x_2}P_2)\\
 &\quad=(-1)^{p_1(n-p_2)}
 N_{P_1}|_{x_1}\times N_{P_2}|_{x_2}
 \times T_{x_1}P_1\times T_{x_2}P_2.
\end{aligned}
\]

（见约定~\ref{conv:8-2-1}(4)。）现对两个浸没
\[
 (\operatorname{ev}_1,\operatorname{ev}_2)\colon U\to U_1\times U_2,
 \qquad
 \mathrm{id}\colon U_1\times U_2\to U_1\times U_2
\]
考虑仓西结构的纤维积
\[
 \left(s\oplus s_1\oplus s_2;
 \mathcal E\oplus N_{P_1}\oplus N_{P_2}
 \longrightarrow
 U_{(\operatorname{ev}_1,\operatorname{ev}_2)}
 \mathop{\times}_{\mathrm{id}}(U_1\times U_2)\right).
\]
它在
$\mathcal M_3^{\mathrm{main}}(L;\beta_0)_{
(\operatorname{ev}_1,\operatorname{ev}_2)}
\mathop{\times}_{f_1\times f_2}(P_1\times P_2)$
上给出仓西结构。对定向而言，由约定~\ref{conv:8-2-1}(4)
得
\begin{equation}
\begin{aligned}
 &\mathcal M_3^{\mathrm{main}}(L;0)_{
 (\operatorname{ev}_1,\operatorname{ev}_2)}
 \mathop{\times}_{f_1\times f_2}(P_1\times P_2)\\
 &\quad=(-1)^{\operatorname{rank}(N_{P_1}\oplus N_{P_2})
 (\dim\mathcal M_3^{\mathrm{main}}(L;0)-\dim(U_1\times U_2))}\\[-2pt]
 &\qquad\qquad{}·(-1)^{p_1(n-p_2)}
 (s\oplus s_1\oplus s_2)^{-1}(0)\\
 &\quad=(-1)^{n(p_1+p_2)+p_1(n-p_2)}
 (s\oplus s_1\oplus s_2)^{-1}(0).
\end{aligned}
\tag{8.6.5}\label{eq:8-6-5}
\end{equation}

现在比较 $(s\oplus s_1\oplus s_2)^{-1}(0)$ 上的定向与
$P_1\cap P_2$ 上的定向。前者的定向由
\[
 \mathcal E_*\times N_{P_1*}\times N_{P_2*}
 \times T_*(s\oplus s_1\oplus s_2)^{-1}(0)
 =T_*\left(
 U_{(\operatorname{ev}_1,\operatorname{ev}_2)}
 \mathop{\times}_{\mathrm{id}}(U_1\times U_2)\right)
\]
定义。下面证明所需结论。

\begin{lemma}{8.6.6}\label{lem:8-6-6}
同构
\[
 \operatorname{ev}_0\colon(s\oplus s_1\oplus s_2)^{-1}(0)
 \longrightarrow P_1\cap P_2
\]
保定向。
\end{lemma}

\begin{proof}[引理~\ref{lem:8-6-6} 的证明]
$P_1\cap P_2$ 和 $(s\oplus s_1\oplus s_2)^{-1}(0)$ 的定向分别由下列
正合列给出：
\[
 0\longrightarrow T_*(P_1\cap P_2)\longrightarrow T_*L
 \longrightarrow N_{P_1}\oplus N_{P_2}\longrightarrow0
\]
与
\[
\begin{aligned}
 0\longrightarrow{}&T_*(s\oplus s_1\oplus s_2)^{-1}(0)
 \longrightarrow
 T_*\left(U_{(\operatorname{ev}_1,\operatorname{ev}_2)}
 \mathop{\times}_{\mathrm{id}}(U_1\times U_2)\right)\\
 &\longrightarrow
 \mathcal E\oplus\pi_1^*N_{P_1}\oplus\pi_2^*N_{P_2}
 \longrightarrow0.
\end{aligned}
\]
考虑把这两个正合列联系起来的下列交换图：

\begingroup
\footnotesize
\[
\begin{array}{ccccc}
T_*(P_1\cap P_2)&\longrightarrow
&\mathcal E_*\oplus T_*L&\longrightarrow
&\mathcal E_*\oplus N_{P_1}\oplus N_{P_2}\\[0.6em]
{\scriptstyle I}\Big\uparrow&&{\scriptstyle\Phi}\Big\uparrow
&&{\scriptstyle\Psi}\Big\uparrow\\[0.6em]
T_*(s\oplus s_1\oplus s_2)^{-1}(0)&\longrightarrow
&\begin{gathered}
T_*\left(U_{(\operatorname{ev}_1,\operatorname{ev}_2)}
\mathop{\times}_{\mathrm{id}}(U_1\times U_2)\right)
\end{gathered}&\longrightarrow
&\begin{gathered}
\mathcal E\oplus\pi_1^*N_{P_1}\oplus\pi_2^*N_{P_2}
\end{gathered}.
\end{array}
\]
\endgroup

顶行在 $\mathcal E_*$ 因子上的映射是恒等映射。
其中的竖直同构如下给出。$I$ 是仓西邻域上评价映射
\[
 \operatorname{ev}_0\colon
 U_{(\operatorname{ev}_1,\operatorname{ev}_2)}
 \mathop{\times}_{\mathrm{id}}(U_1\times U_2)\longrightarrow L
\]
的微分在 $T_*(s\oplus s_1\oplus s_2)^{-1}(0)$ 上的限制。换言之，
\[
 I=T_*\operatorname{ev}_0\colon
 T_*(s\oplus s_1\oplus s_2)^{-1}(0)
 \longrightarrow T_*(P_1\cap P_2),
\]
并且可显式写成
\[
 I((0,v),v,v)=v.
\]
映射 $\Phi$ 和 $\Psi$ 定义为
\[
 \Phi((a,b),v_1,v_2)
 =\bigl(a,\operatorname{pr}_1(v_1)+\operatorname{pr}_2(v_2)
 +\operatorname{pr}_3(b)\bigr),
 \qquad v_1=a+b,\quad v_2=b,
\]
以及
\[
 \Psi(u,v_1,v_2)=(u,v_1,v_2).
\]
这里 $\operatorname{pr}_i$ 是直和分解
\[
 T_*L=N_{P_1}\oplus N_{P_2}\oplus T_*(P_1\cap P_2)
\]
的第 $i$ 个因子投影。%
\translatornote{源书此处把第三个直和因子写成空间
$P_1\cap P_2$，但左端与前两个因子均为切向量空间，且
$\operatorname{pr}_3(b)$ 也必须是向量；故恢复为
$T_*(P_1\cap P_2)$。}
由于同构 $\Phi$ 和 $\Psi$ 保定向，追图论证表明同构 $I$
也保定向。引理得证。
\end{proof}

因此，结合式~\eqref{eq:8-6-5} 与引理~\ref{lem:8-6-6}，
就完成了命题~\ref{prop:8-6-3} 的证明。
\end{proof}
